\chapter{Cenni di storia dell'anatomia}

% ---------------------------------------------------------------------------
\section{Definizione di anatomia}

Il termine \textbf{anatomia} deriva dal greco \textit{anatomè} (``dissezione''), formato da \textit{anà} (``attraverso'') e \textit{tèmno} (``tagliare''). È la branca della biologia che studia la \textbf{forma e la struttura} degli organismi.

L'anatomia si suddivide ulteriormente in:
\begin{itemize}
  \item \textbf{normale} e \textbf{patologica};
  \item artistica, superficiale, di comparazione anatomica, esperienziale.
\end{itemize}

L'anatomia e la fisiologia studiano insieme la forma e il funzionamento del corpo.

% ---------------------------------------------------------------------------
\section{Breve storia dell'anatomia}

\subsection{L'Antico Egitto}

Gli egizi seguivano nelle loro trattazioni mediche un ordine pressoché geografico, partendo dall'alto --- cioè dal capo --- e proseguendo descrittivamente fino ai piedi. Le loro conoscenze anatomiche sono documentate da diversi papiri:

\begin{itemize}
  \item \textbf{1850 a.C.} --- Papiro di Kahun: trattato di ginecologia.
  \item \textbf{1550 a.C.} --- Papiro di Ebers: anatomia e patologie; Papiro di Hearst: apparato urinario, sangue, denti.
  \item Documenti su chirurgia delle fratture e medicamenti.
\end{itemize}

\noindent Tra i testi sacri, il \textit{Libro dei Morti di Hunefer} (British Museum, Londra) descrive la cerimonia di apertura della bocca, testimoniando una conoscenza pratica del corpo legata al rito funerario.

\subsection{La Grecia antica e il Medioevo}

I primi esperimenti anatomici sistematici nascono in Grecia. La \textbf{dissezione dei cadaveri} avviene già nel XIII secolo: dapprima come esame autoptico del corpo di chi era morto in circostanze dubbie, poi a scopo didattico, per conoscere meglio gli organi e l'architettura del corpo umano.

La prima sede di rilievo fu \textbf{Bologna}: dopo il testo di chirurgia del chierico Guglielmo da Saliceto (1270), che si basava sull'osservazione diretta di corpi dissezionati, viene redatta l'\textbf{\textit{Anathomia Mundini}} del 1316, di Mondino dei Liuzzi --- il trattato più importante di anatomia medievale.

In seguito, la \textbf{scuola anatomica di Padova} produce l'opera fondamentale di \textbf{Andrea Vesalius}, \textit{De humani corporis fabrica}, pubblicata nel \textbf{1543}: nasce così l'Anatomia Moderna, nel periodo del Rinascimento, basata sull'osservazione diretta su cadaveri.

Nel Medioevo la medicina era prevalentemente:
\begin{itemize}
  \item pratica e diretta;
  \item basata su cure con stregonerie, incantesimi e terapie a base d'erbe;
  \item associata alla nascita dei primi ospedali, ospizi e monasteri.
\end{itemize}

\noindent Una curiosità del Medioevo è il cosiddetto \textbf{\textit{wound man}}: una figura anatomica impiegata come importante guida per i chirurghi (o i barbieri-chirurghi) rispetto alle potenziali ferite. Lo scopo era quello di rivelare le potenziali cure e i trattamenti disponibili al tempo (es.\ manoscritto della fine del XV secolo, Bayerische Staatsbibliothek Cgm 597).

Gli strumenti chirurgici tipici del XV secolo comprendevano: forbici, forbice-pinza, uncino, forcella, sonda, bisturi, aghi, filo e sega.

\subsection{Mondino de' Liuzzi}

\textbf{Mondino de' Liuzzi} è stato uno dei più grandi anatomisti italiani (attivo intorno al 1300). Utilizzò cadaveri di criminali per le sue dissezioni, lottando contro il tempo a causa del rapido deterioramento dei corpi. Nel \textbf{1316} pubblicò il suo manuale di anatomia, noto come \textit{Anothomia}.

\subsection{Il Rinascimento: Leonardo da Vinci}

\textbf{Leonardo da Vinci} rappresenta una svolta fondamentale nella storia dell'anatomia. L'occhio dell'artista vede linee, volumi e la perfetta composizione delle proporzioni. Leonardo stesso scrisse:

\begin{quote}
\textit{``\ldots questa mia figurazione del corpo umano ti sarà dimostrata non altrementi che se tu avessi l'omo naturale innanzi.''}\\
\hfill --- Leonardo da Vinci
\end{quote}

La vera innovazione nei suoi studi è la \textbf{ricerca della funzione} degli organi e degli apparati da lui esaminati. Leonardo applica costantemente lo studio fisiologico a quello anatomico; caratteristica della sua ``ricerca della funzione'' sono le \textbf{vedute esplose}, che scompongono nello spazio i particolari per meglio comprenderli. Questo contributo è definibile come un impulso all'\textbf{anatomo-fisiologia}.

Leonardo eseguì dissezioni in prima persona: famosa è la narrazione autoptica del ``vecchio di cent'anni'', morto placidamente all'Ospedale di Santa Maria Nuova di Firenze, del quale descrisse la causa della morte dopo averlo sezionato.

\noindent I disegni anatomici di Leonardo --- relativi al cranio, alla colonna vertebrale, agli arti, agli organi interni e al feto --- furono realizzati indicativamente tra il 1480 e il 1517.

\subsection{Guido da Vigevano}

\textbf{Guido da Vigevano} è un'altra figura di rilievo. Il suo trattato \textit{Anothomia designata per figuras} (1345) costituisce uno dei primi esempi di anatomia illustrata. La sua opera è stata posta a confronto con quella di Leonardo da Vinci in una mostra tenutasi alla Biblioteca Ambrosiana di Milano (6 febbraio 2020), che ha presentato le riproduzioni di 18 figure (in 16 tavole) del trattato, accostate a disegni anatomici di Leonardo.

\subsection{L'epoca moderna: dal 1800 ad oggi}

Nel \textbf{1800} si registrano episodi di trafugamento di cadaveri, sezionamento di soggetti vivi e abusi connessi alla pratica anatomica. Questo periodo segna anche l'\textbf{inizio della storia dell'Anatomia Patologica}.

Si afferma progressivamente lo \textbf{studio della Morfologia}, che studia la forma delle strutture e i rapporti tra i diversi organi.

Oggi l'anatomia studia la struttura a \textbf{molteplici livelli di osservazione}:
\begin{itemize}
  \item tessuti;
  \item cellule;
  \item molecole.
\end{itemize}

Il corpo umano consiste di un sistema biologico composto da organi e apparati, costituiti da tessuto biologico, il quale è a sua volta composto da cellule biologiche e tessuto connettivo.

Un punto di riferimento storico imprescindibile è l'\textbf{Anatomia del Gray} (\textit{Anatomy of the Human Body}, di Henry Gray), che rappresenta --- non solo per gli studenti di medicina e per gli specializzandi, ma anche per i medici di medicina generale e per gli specialisti --- un punto di riferimento per l'apprendimento dell'anatomia, disciplina considerata la base della pratica clinica.

% ---------------------------------------------------------------------------
\section{Anatomia microscopica: citologia e istologia}

\subsection{Anatomia microscopica}

L'\textbf{anatomia microscopica} studia quelle strutture che \textbf{non} possono essere viste senza l'ingrandimento microscopico. Si articola in:

\begin{itemize}
  \item \textbf{Citologia}: analizza i singoli componenti cellulari.
  \item \textbf{Istologia}: esamina i tessuti, ovvero l'insieme di cellule e prodotti cellulari che, cooperando tra loro, determinano lo svolgersi di una o più specifiche funzioni.
\end{itemize}

Tutte le cellule del corpo possono essere suddivise in \textbf{4 tipi tessutali}:
\begin{enumerate}
  \item Epiteli
  \item Tessuti connettivi
  \item Tessuto muscolare
  \item Tessuto nervoso
\end{enumerate}

\subsection{Istologia}

L'istologia è una branca della biologia applicata alla medicina, con un ruolo fondamentale anche in Anatomia Patologica.

\begin{itemize}
  \item Il barone \textbf{Albrecht von Haller} (1708--1777) suggerì per primo un'organizzazione strutturale microscopica degli esseri viventi.
  \item L'\textbf{istologia moderna} nacque nel XIX secolo con la \textbf{teoria cellulare}.
\end{itemize}

La metodologia si basa sul microscopio, che permette l'osservazione diretta dei tessuti: questi vengono tagliati in strisce sottilissime, trattati per prevenire la decomposizione e infine colorati per rendere visibili le strutture.

% ---------------------------------------------------------------------------
\section{Organizzazione del corpo umano}

Lo studio anatomico moderno si svolge a due livelli principali:

\begin{itemize}
  \item \textbf{Organizzazione macroscopica}: indagata tramite tecniche di imaging (radiografia, TAC, RMN, ecografia, PET), che permettono di esplorare l'interno del corpo senza dissezione.
  \item \textbf{Organizzazione microscopica}: indagata tramite istologia e citologia.
\end{itemize}

Un esempio contemporaneo di anatomia applicata è la mostra \textit{Real Bodies} (Milano), in cui cadaveri plastinati vengono esposti al pubblico per illustrare l'organizzazione interna del corpo umano.

% ---------------------------------------------------------------------------
\section{L'anatomia e il futuro}

L'anatomia è una disciplina in continua evoluzione: dall'osservazione macroscopica dei cadaveri si è giunti all'imaging digitale tridimensionale, all'integrazione con la biomeccanica e alle prospettive offerte dalla bioingegneria (protesi neurali, arti artificiali). L'anatomia e la fisiologia rimangono le basi irrinunciabili della comprensione del corpo umano e della pratica clinica.
